\chapter*{Abstract}
\addcontentsline{toc}{chapter}{Abstract}

Large Language Models (LLMs) have showed potential powers in natural language understanding and generation, yet their direct application in educational settings presents significant challenges. General-purpose models often produce hallucinated content, fail to adapt to individual student knowledge levels, and lack awareness of course-specific materials, limiting their effectiveness as pedagogical tools. This thesis addresses these limitations by proposing \textbf{TripartiteRAFT}, a novel framework that integrates Retrieval-Augmented Generation (RAG), domain-specific fine-tuning, and personalized response generation into a one pipeline for educational question answering.

The proposed system operates in three complementary stages. First, a hybrid retrieval mechanism combines dense vector search with prerequisite concept gathering to identify both directly relevant content and foundational knowledge required by the student. Second, the prerequisites and symbols associated with each piece of retrieved content are extracted. Finally, these prerequisites are classified according to the learner's existing knowledge background, as represented by the learner model. The resulting information is then used to construct a complete prompt, which is subsequently passed to the fine-tuned model for response generation.

To support this work, two datasets were constructed: a collection of real student questions sourced from forum interactions, and a handcrafted dataset of question--answer pairs annotated with student knowledge backgrounds. The system was evaluated using multiple methodologies, including LLM-as-a-judge metrics for contextual relevancy, answer relevancy, and faithfulness, as well as human expert assessment. Experiments with various embedding models were also conducted to analyze the impact of embedding dimensionality on retrieval quality.

The results demonstrate that TripartiteRAFT significantly outperforms baseline LLM responses in terms of factual accuracy, contextual grounding, and pedagogical relevance. The fine-tuned model exhibits a markedly reduced propensity for hallucinations and shows improved ability to incorporate course-specific terminology and conceptual frameworks.

This work contributes a complete, deployable pipeline for personalized educational question answering, with potential for integration into platforms such as ALeA. The findings underscore the promise of combining retrieval augmentation with fine-tuning as a pathway toward safer, more effective, and more equitable AI-assisted learning environments.